% Part: propositional-logic
% Chapter: propositional-logic 
% Section: soundness

\documentclass[../../../include/open-logic-section]{subfiles}

\begin{document}

\olfileid{pl}{prp}{snd}

\olsection{પ્રતિજ્ઞાપન તર્કશાસ્ત્રની યથાર્થતા}

\begin{thm}[યથાર્થતા]
\ollabel{thm:soundness}
જો $\Gamma\Proves!A$ હોય, તો $\Gamma \Entails !A$.
\end{thm}

\begin{proof}
નિષ્પન્ન થતી પ્રતિજ્ઞાપનાઓ પર આગમન કરીએ.
જો $!A$ સ્વયંસિદ્ધ હોય, તો $\Entails
!A$, એટલે $\Gamma \Entails!A$ પણ થાય.
તે જ રીતે, જો $!A \in \Gamma$ હોય, તો $\Gamma \Entails!A$.
આગમનના પગલામાં ધારો કે $!A$ એ $!C$ અને $!C\lif!A$માંથી
\emph{મોડસ પોનેન્સ} વડે મળ્યું હોય.
આગમનની ધારણા મુજબ $!C\lif !A$ અને $!C$ માટે પ્રમેય સાચું છે.
\olref[pl][syn][sem]{prop:semanticalfacts}નો બીજો મુદ્દો
ઇચ્છિત પરિણામ આપે છે.
\sourcecorrection{OLMD-345}{મૂળ અહીં અર્થવિજ્ઞાનના પરિણામનો ત્રીજો મુદ્દો કહે છે. સ્થિર મૂળમાં મોડસ પોનેન્સને અનુરૂપ સત્ય-સંરક્ષણ બીજો મુદ્દો છે; ત્રીજો મુદ્દો સાંત ઉપગણોની સંતોષ્યતા અંગે છે. ક્રમસંખ્યા અને સ્રોતના સાચા ભાગ-અધ્યાય સુધીનો સંદર્ભ સુધાર્યા છે.}
\end{proof}

\begin{cor}
જો $\Gamma$ સંતોષ્ય હોય, તો $\Gamma$ સુસંગત છે.
આથી પ્રતિજ્ઞાપન તર્કશાસ્ત્ર સુસંગત છે.
\end{cor}

\end{document}

